% !TeX root = ../教材题库_全部条目_转写工作稿.tex
% Source: XJ-B1-C06-HW-P214-Y0496
% 第6章 / 6.1 获取数据的途径与统计的一些基本概念 / 教材 p.214 / PDF 第221页
% OCR未校对：公式与文字均可能有误；未完成数学TeX转写；禁止直接发给学生
% 原题及相邻原解截图（核对用）：../题目原图/XJ-B1-C06-HW-P214-Y0496-01.webp
% 本文件只在转写完成后作为题干使用；OCR 条目尚不是可靠题干。
\begin{bankitem}{习题6.1 第2题}{第6章 / 6.1 获取数据的途径与统计的一些基本概念 / 教材 p.214 / PDF 第221页}{XJ-B1-C06-HW-P214-Y0496}{OCR未校对：公式与文字均可能有误；未完成数学TeX转写；禁止直接发给学生}
\begin{ocrdraft}
[教材p.214 / PDF第221页]
\par
2．
\par
为调查参加运动会的1000名运动员的年龄情况，从中抽查了100名运动员
\par
的年龄，下列说法正确的是（）
\par
A.
\par
1000名运动员是总体
\par
B.每个运动员是个体
\par
C.抽取的100名运动员是样本
\par
D.样本容量是100
\end{ocrdraft}
\end{bankitem}
